% Part: normal-modal-logic
% Chapter: completeness
% Section: complete-consistent-sets

\documentclass[../../../include/open-logic-section]{subfiles}

\begin{document}

\olfileid[pt]{nml}{com}{ccs}

\olsection{Conjuntos $\Sigma$-consistentes completos}

Suponha que $\Sigma$ seja um conjunto de !!{formula}s modais---pense
nelas como os axiomas ou princípios definidores de uma lógica modal
normal. Um conjunto $\Gamma$ é $\Sigma$-consistente se e somente se
$\Gamma \Proves/[\Sigma] \lfalse$, O.s.\@, se não existe
!!{derivation} de~$!A_1 \lif (!A_2 \lif \cdots (!A_n \lif
\lfalse)\dots)$ a partir de $\Sigma$, onde cada $!A_i \in \Gamma$.
Construiremos um modelo ``canônico'' no qual cada mundo é tomado como
um tipo especial de conjunto $\Sigma$-consistente: um conjunto que não
é apenas $\Sigma$-consistente, mas maximamente consistente, no sentido
de que determina o valor de verdade de toda !!{formula} modal: para
todo $!A$, ou $!A \in \Gamma$, ou $\lnot !A \in \Gamma$:

\begin{defn}
  Um conjunto $\Gamma$ é \emph{$\Sigma$-consistente completo} se e
  somente se é $\Sigma$-consistente e, para todo~$!A$, ou
  $!A \in \Gamma$, ou $\lnot !A \in \Gamma$.
\end{defn}

Conjuntos $\Sigma$-consistentes completos $\Gamma$ têm várias
propriedades úteis. Em primeiro lugar, são dedutivamente fechados, O.s.\@, se $\Gamma \Proves[\Sigma] !A$, então $!A \in \Gamma$. Isso
significa, em particular, que toda instância de !!a{formula}
$!A \in \Sigma$ também está~$\in \Gamma$. Além disso, a pertinência
a $\Gamma$ reflete as condições de verdade dos conectivos
proposicionais. Isso será importante quando definirmos o ``modelo
canônico''.

\begin{prop}\ollabel{prop:ccs-properties}
  Suponha que $\Gamma$ seja $\Sigma$-consistente completo. Então:
  \begin{enumerate}
  \item \ollabel{prop:ccs-closed}%
    $\Gamma$ é dedutivamente fechado em~$\Sigma$.
  \item \ollabel{prop:ccs-sigma}%
    $\Sigma \subseteq \Gamma$.
  \tagitem{prvFalse}{\ollabel{prop:ccs-lfalse}%
    $\lfalse \notin \Gamma$}{}
  \tagitem{prvTrue}{\ollabel{prop:ccs-ltrue}%
    $\ltrue \in \Gamma$}{}
  \item \ollabel{prop:ccs-lnot}%
    $\lnot!A \in \Gamma$ se e somente se $!A \notin
    \Gamma$.
  \tagitem{prvAnd}{\ollabel{prop:ccs-land}%
    $!A \land !B \in \Gamma$ se e somente se $!A \in \Gamma$ e $!B \in \Gamma$}{}
  \tagitem{prvOr}{\ollabel{prop:ccs-lor}%
    $!A \lor !B \in \Gamma$ se e somente se $!A \in \Gamma$ ou $!B \in \Gamma$}{}
  \tagitem{prvIf}{\ollabel{prop:ccs-lif}%
    $!A \lif !B \in \Gamma$ se e somente se $!A \notin \Gamma$ ou $!B \in \Gamma$}{}
  \tagitem{prvIff}{\ollabel{prop:ccs-liff}%
    $!A \liff !B \in \Gamma$ se e somente se ou $!A \in \Gamma$ e $!B \in
    \Gamma$, ou $!A \notin \Gamma$ e $!B \notin \Gamma$}{}
  \end{enumerate}
\end{prop}

\begin{proof}
  \begin{enumerate}
  \item Suponha que $\Gamma \Proves[\Sigma] !A$, mas $!A \notin
    \Gamma$. Então, como $\Gamma$ é $\Sigma$-consistente completo,
    $\lnot!A \in \Gamma$. Isso tornaria $\Gamma$ inconsistente, pois
    $!A, \lnot !A \Proves[\Sigma] \lfalse$.
    
  \item Se $!A \in \Sigma$, então $\Gamma \Proves[\Sigma] !A$, e $!A
    \in \Gamma$ pelo fechamento dedutivo, O.s.\@, pelo
    caso~\olref{prop:ccs-closed}.
    
  \tagitem{prvFalse}{Se $\lfalse \in \Gamma$, então $\Gamma
    \Proves[\Sigma] \lfalse$, de modo que $\Gamma$ seria
    $\Sigma$-inconsistente.}{}
  
  \tagitem{prvTrue}{$\Gamma \Proves[\Sigma] \ltrue$, logo $\ltrue \in
    \Gamma$ pelo fechamento dedutivo, O.s.\@, pelo
    caso~\olref{prop:ccs-closed}.}{}

\item Se $\lnot!A \in \Gamma$, então, pela consistência, $!A \notin
    \Gamma$; e, se $!A \notin \Gamma$, então $\lnot !A \in \Gamma$, pois
    $\Gamma$ é $\Sigma$-consistente completo.
    
  \tagitem{prvAnd}{\iftag{probAnd}{Exercício.}{Suponha que $!A \land !B \in
      \Gamma$. Como $(!A \land !B) \lif !A$ é uma instância
      tautológica, $!A \in \Gamma$ pelo fechamento dedutivo, O.s.\@,
      pelo caso~\olref{prop:ccs-closed}. O mesmo vale para $!B \in
      \Gamma$. Por outro lado, suponha tanto $!A \in \Gamma$ quanto
      $!B \in \Gamma$. Então, o fechamento dedutivo implica
      $(!A \land !B) \in \Gamma$, pois
      $!A \lif (!B \lif (!A \land !B))$ é uma instância tautológica.}}{}

  \tagitem{prvOr}{\iftag{probOr}{Exercício.}{Suponha que $!A \lor !B \in
      \Gamma$, e que $!A \notin \Gamma$ e $!B \notin \Gamma$. Como
      $\Gamma$ é $\Sigma$-consistente completo, $\lnot!A \in \Gamma$
      e $\lnot !B \in \Gamma$. Então $\lnot(!A \lor !B) \in \Gamma$,
      pois $\lnot !A \lif (\lnot !B \lif \lnot (!A \lor !B))$ é uma
      instância tautológica. Isso significaria que $\Gamma$ é
      $\Sigma$-inconsistente, uma contradição.}}{}
  
  \tagitem{prvIf}{\iftag{probIf}{Exercício.}{Suponha que $!A \lif !B \in
      \Gamma$ e $!A \in \Gamma$; então $\Gamma \Proves[\Sigma] !B$,
      de onde $!B \in \Gamma$ pelo fechamento dedutivo. Reciprocamente,
      se $!A \lif !B \notin \Gamma$, então, como $\Gamma$ é
      $\Sigma$-consistente completo, $\lnot (!A \lif !B) \in \Gamma$.
      Como $\lnot(!A \lif !B) \lif !A$ é uma instância tautológica,
      $!A \in \Gamma$ pelo fechamento dedutivo. Como
      $\lnot(!A \lif !B) \lif \lnot !B$ é uma instância tautológica,
      $\lnot !B \in \Gamma$. Então $!B \notin \Gamma$, pois $\Gamma$
      é $\Sigma$-consistente.}}{}
  
  \tagitem{prvIff}{\iftag{probIff}{Exercício.}{Suponha que $!A \liff !B \in
      \Gamma$. Se $!A \in \Gamma$, então $!B \in \Gamma$, pois $(!A
      \liff !B) \lif (!A \lif !B)$ é uma instância tautológica. De
      modo semelhante, se $!B \in \Gamma$, então $!A \in \Gamma$.
      Portanto, ou tanto $!A \in \Gamma$ quanto $!B \in \Gamma$, ou
      nem $!A \in \Gamma$ nem $!B \in \Gamma$.
      
      Reciprocamente, suponha que $!A \lif !B \notin \Gamma$. Como
      $\Gamma$ é $\Sigma$-consistente completo,
      $\lnot (!A \liff !B) \in \Gamma$. Como
      $\lnot(!A \liff !B) \lif (!A \lif \lnot !B)$ é uma instância
      tautológica, se $!A \in \Gamma$, então $\lnot !B \in \Gamma$;
      e, como $\Gamma$ é $\Sigma$-consistente, $!B \notin \Gamma$. De
      modo semelhante, se $!B \in \Gamma$, então $!A \notin \Gamma$.
      Portanto, não ocorre nem $!A \in \Gamma$ e $!B \in \Gamma$, nem
      $!A \notin \Gamma$ e $!B \notin \Gamma$.}}{}
  \end{enumerate}
\end{proof}

\begin{probtag}{probAnd,probOr,probIf,probIff}
  Complete a demonstração de \olref[nml][com][ccs]{prop:ccs-properties}.
\end{probtag}

\end{document}
